\section{Conclusion and the second stage}\label{sec:conclusion}

\subsection{What this stage shows}
This report set out to measure the band within which Indian stock futures are priced relative to their stocks, to ask
whether costs and frictions set that band, and to ask whether the F\&O ban restricts trading as intended. The answers
are as follows.

The basis of Indian stock futures moves in a wide, asymmetric band, from about 35\bp\ below each stock's typical level
to about 24\bp\ above it. The band is nonlinear in the way theory predicts, it is stable over time apart from the turmoil
of 2020, and its lower side reaches further, as expected when the trade that corrects cheap futures needs a short sale.
It is much wider than the intraday bands found in earlier Indian studies, because an arbitrageur who carries a position
overnight pays the taxes on delivery. Daily prices contain enough noise that the speed of reversion should be read with
care; where the band lies is the more reliable finding.

Costs shape the band, but loosely. The full cost of a new cash-and-carry position lies in the outer tail of the basis,
and the much smaller cost of rolling an existing position helps explain why the estimated edge lies inside it. The upper
part of the basis rose after the 2024 and 2026 tax increases, as predicted and as the test fixed in advance confirms, but
by more than the cost changes and by amounts that are not unusual on other dates. Across stocks, the band is wider where
trading is harder.

The ban binds hard on trading. Futures trading falls by about 87\% on ban days, cash trading moves toward delivery, and
once violations were judged on intraday positions every futures trade during a ban closed an existing position.
Enforcement therefore matters as much as the rule itself. Traders do not hold back before the trigger, and the 2025
rules cut the number of ban days by about three quarters, mainly by measuring open interest in delta terms.

\subsection{Limitations}
The data are daily, so trading within the day is outside the scope of this report, and noise in daily prices inflates
measured reversion. The positions of arbitrageurs are not observed, so the model's prediction that the band moves with
the existing book can only be checked indirectly. The tax increases are three events, one of which coincides with a
change in fund taxation, so the test of costs has little power against shocks common to all stocks. The period after the
2025 reform is short.

\subsection{The second stage}
The first stage establishes the model, the test plan, the data, the band and the ban's effect on trading. The second
stage uses these to ask what the limits do to prices:
\begin{itemize}[leftmargin=1.4em,itemsep=2pt]
\item \textbf{The ban and the basis.} Whether correction of the basis becomes one-sided in a ban
(\cref{prop:onesided}, H7, a primary hypothesis), how the level of the basis changes when a ban starts and whether this
depends on the composition of trading (\cref{prop:level}), the 95\% trigger as a discontinuity in prices, and placebo
bans at false dates.
\item \textbf{Position limits before a ban.} Whether the basis becomes richer as utilisation approaches the trigger, and
what the 2025 formula and the 2020 halving and restoration of limits did to prices.
\item \textbf{Arbitrage capital.} Whether flows into arbitrage funds narrow the band.
\item \textbf{Short-sale constraints across the market.} Whether entry into the F\&O segment, which brings the right to sell
short, speeds the response of prices to news (H6, a primary hypothesis), whether anomalies are stronger where shorting
is hardest, and a measure of short-sale constraints for every listed stock.
\item \textbf{Policy.} What the results imply for the design and enforcement of position limits.
\end{itemize}
